\subsection*{Problems}

\textbf{11.1.} Suppose r< s. Show that if E[\vertX\verts] is finite, then E[\vertX\vertr] is also finite.

\textbf{11.2.} (Cantelli inequality) Let X be a random variable with finite variance, and let

a \geq 0. Show that

P(X - E[X]\geq a) \leq Va r(X)

Va r(X) + a2 .

\textbf{11.3.} Let X1,X 2,...,X n be iid. random variables uniformly distributed between

0 and 1. Show that the sequence of random variables

Yn \coloneqq X1 + X2 +\cdot\cdot\cdot+ Xn

X2

1 + X2

2 +\cdot\cdot\cdot+ X2n

converges to a constant in probability as n \to\infty .

\textbf{11.4.} Let (Xn)n\geq1 be iid. random variables with pdf f( x) = 2x- 31[1,\infty )(x).

Apply Khinchin's weak law to show that (X 1 + X2 +\cdot\cdot\cdot+ Xn)/n

P

\to 2.

\textbf{11.5.} Suppose that Xk ,f o r k = 1, 2, 3,... , are independent random variables with

zero mean and variance V ar (Xk) \leq c

\sqrt

k for some positive constant c. Prove that

Xn/n converges almost surely to zero as n \to\infty .

\textbf{11.6.} Suppose (Xk)\infty

k=1 is a sequence of iid. random variables with zero mean, and

suppose the random variables are uniformly bounded, ie., there is a constant c such

that P( \vertXk\vert <c ) = 1 for all kLet Sn = X1 +\cdot\cdot\cdot+ Xn. Prove that Sn/n3/4 as.

-\to 0

as n \to\infty .

(Hint: Show that E[S6

n] is upper bounded by cn3 for some constant c.)

\textbf{11.7.} Let (Xi)\infty

i=1 be a sequence of random variables that may be correlated.

Suppose E [Xi]= \mu and V ar(Xi) \leq K for all i. Show that if there exists a sequence

of real numbers (a\tau)\tau\geq1 such that (i) a\tau \in[ 0, 1] for all \tau, (ii) a\tau decreases to 0 as

\tau increases, and (iii) Cov (Xi,Xi+\tau) \leq a\tau

\sqrt

Va r(Xi) Va r(Xi+\tau) for all i and \tau> 0,

then (Xi)\infty

i=1 converges in probability to \mu.

\textbf{11.8.} Define a first-order autocorrelation process by X0 = 0 and Xi = \rhoXi- 1 + Ni

for i \geq 1, where \rho is a constant with \vert\rho\vert < 1 and Ni is a Gaussian random variable

N( 0,\sigma 2)Assume that the random variables Ni 's are independent. Show that

(X1 + X2 +\cdot\cdot\cdot+ Xi)/i converges to 0 in probability as i \to\infty .

\textbf{11.9.} Consider the experiment of throwing k distinct balls into n distinct boxes.

Assume the locations of the balls are independent and uniformly chosen among

{1, 2,...,n }We consider an asymptotic scenario in which k and n increase

simultaneously with k/n \to a, for some positive constant aLet Xn denote the

number of empty boxes when the number of boxes is n Prove that X n/n \to e-a in

quadratic mean, and hence in probability as n \to\infty .

\textbf{11.10.} Prove Glivenko-Cantelli theorem under a simplifying assumption that the

limit cdf F(x) is continuous.
